\section{Paid asynchronous jobs and on-chain
coordination}\label{paid-asynchronous-jobs-and-on-chain-coordination}

B transformation is the primary commercial product: author-directed
cleanup that preserves meaning and target tone while pursuing the
required detector objectives. Customers pay for access to competing
private models, accepted-revision performance, convenience and qualified
attested execution. Private weights are not released on a calendar;
miners can retain improvements while offering newly benchmarked versions.
A remains an open detector resource and a secondary hosted service for
customers who prefer an API to running it themselves. The market determines
whether A hosting is cheaper; the protocol mandates no price ratio.

Miners earn network emissions for independently evaluated benchmark
utility and may separately earn customer fees. These are different
revenue sources. Customers can pay for B without receiving its weights;
benchmark version commitments and attested receipts support that service
claim. Hardware attestation supplies no exclusivity, guaranteed premium
or permanent advantage against detector providers studying served outputs.
Reader outcomes, repeat purchases and margins after confidential-compute
costs must test the business thesis.

The proposed subnet settlement transfers the quoted fee to the serving
miner, subject to its disclosed fees and payment rules. It does not
automatically pay a royalty to Slop Ninja Research or make customer
revenue an emission multiplier. Any future application subscription or
gateway fee needs explicit customer terms; none is assumed in the
protocol's economics. Alpha-denominated demand alone establishes no
token-price or investor-return guarantee.

Offers identify ordinary or attested execution, model/version evidence
and the accepted runtime policy. Sensitive B jobs require attestation
under Section~\ref{attested-customer-inference}; a customer must explicitly
choose ordinary hosting to grant the operator plaintext access. Public A
and voluntarily published B models can run locally, without a subnet
job or fee. No public copy of the best private B model is promised.

B jobs pursue author matching and detector evasion within the selected
execution boundary. A confidential job has no verified Pangram result
unless the customer separately authorizes external disclosure; its
objectives do not expand recipient access. Authorized benchmark rounds
provide the routine external measurements under
Section~\ref{pangram-cadence-and-parity}.
The proposed job contract uses Subtensor's EVM to coordinate offers,
escrow, commitments, and settlement. Inference prices are denominated in
Slop Ninja's subnet alpha; execution stays off-chain.
Subtensor EVM guide~\cite{evm}.

\subsection{Market pricing in alpha}\label{market-pricing-in-alpha}

Miners publish signed offers for bounded jobs. Each offer specifies a
total alpha price, task, service version and execution mode, input/output limits,
delivery deadline, available capacity, quote expiry, and settlement
terms. Attested sensitive-job offers also bind the padded workload bucket
and fixed response-release slot; charge the reserved job price without
publishing content-dependent token usage. Customers select among compatible offers using price, benchmark
performance, and delivery time. A higher-quality model or faster service
may command a different price. There need not be one clearing price for
every kind of inference.

Miners may revise their unreserved offers as demand, queue length,
compute costs, and competition change. The subnet operator sets no base
price, utilization target, or global price-adjustment coefficient.
Miners choose their own pricing policies; customers decide whether the
offers are worth buying. Alpha is the unit of account and proposed
settlement asset. The inference market determines the number of alpha
units charged independently of the alpha/TAO exchange market.

An atomic reservation consumes offered capacity and locks the quoted
alpha amount and job terms. Later price changes apply only to later
reservations. An expired or exhausted offer cannot be filled; repricing
requires a new customer choice. The reservation therefore gives both
parties a known price through execution and settlement, even while
other offers change.

Quote discovery uses public service descriptions and the size/deadline
information the customer chooses to reveal. Source text, reference
writing, and the full brief remain encrypted to the selected instance
or, for explicitly ordinary hosting, its operator-controlled key.
Obtaining competing offers does not distribute the plaintext to
prospective bidders. If the selected miner cannot serve the bounded
request, the job follows its rejection/refund terms; it cannot silently
raise the price after decrypting the input.

\subsection{Reservation and delivery}\label{reservation-and-delivery}

Subtensor's staking-v2 interface exposes allowances and stake transfers,
and contracts can hold stake through their mapped coldkeys. The proposed
escrow adapter would use those operations to custody unlocked,
transferable alpha on the same subnet. An allowance alone does not fund
a reservation. Alpha remains associated with a subnet and hotkey;
settlement must account for that position and its native \(10^{-9}\)
precision. These primitives do not supply a finished job-escrow
implementation~\cite{stakingv2,evmstake}.

An offer distinguishes the service price from network and transfer fees
and names the fee payer. Funding succeeds only after the credited alpha
principal satisfies the quote. Release and refund must reconcile actual
balances, transfer permissions, and fees; incidental stake accrual must
remain separate from job principal. The adapter and its accounting
policy require validation before paid launch. Neither an alpha price
feed nor an assumed ERC-20 interface substitutes for that work.

The proposed lifecycle is:

\begin{enumerate}
\def\labelenumi{\arabic{enumi}.}
\tightlist
\item
  \textbf{Offer and reserve.} The customer selects a miner's signed
  offer, pins its execution mode, version and job terms, and escrows the
  quoted alpha amount for a bounded input/output. Attested jobs verify
  fresh instance evidence and its job key before any input disclosure;
  failed verification permits cancellation/refund under the offer.
\item
  \textbf{Encrypt and submit.} The customer posts the salted input
  commitment and encrypted-delivery hash, then delivers the envelope to
  the selected recipient key, including the attested instance binding
  where required.
\item
  \textbf{Accept.} The selected runtime decrypts, checks the opening and limits,
  and accepts before the acceptance deadline. In attested mode this
  occurs inside the protected application, without operator access. Failure to accept permits
  a refund.
\item
  \textbf{Execute and deliver.} The selected runtime runs the private service,
  commits the result, and makes its customer-encrypted envelope
  available before the delivery deadline. Attested execution also returns
  the bound receipt from its protected signing key; the host only relays
  the ciphertext and permitted settlement metadata.
\item
  \textbf{Acknowledge or expire.} A customer acknowledgment releases
  payment. If there is no timely acceptance or delivery claim, the
  customer may claim a refund. A
  delivery claim without acknowledgment follows the bounded timeout
  policy below.
\end{enumerate}

Bind acknowledgments and state transitions to chain, contract, job ID,
role, nonce, and assignment generation. Bind the alpha asset's subnet ID
and amount in native units as well. Authenticate the association
between a Bittensor hotkey and its settlement identity. A hash or
truncated address is not an authorization method. Enforce exactly one
terminal settlement and explicit finality rules. A customer, miner, or
keeper must submit timeout transactions; deadlines do not execute
themselves.

Customers may deliver ciphertext through authenticated asynchronous
endpoints and redundant encrypted storage. Small envelopes could instead
travel through on-chain calldata or events, so miners can discover and
complete jobs entirely through chain messaging. Payload-size, gas, and
retention measurements must precede that choice. Storing ciphertext
on-chain does not hide metadata or provide permanent secrecy after key
compromise.

\subsection{Payment without disclosing the
job}\label{payment-without-disclosing-the-job}

A customer can consume a useful result and refuse acknowledgment. A
miner can claim delivery of useless or undecryptable bytes. Public
commitments alone do not distinguish those cases. Validators who cannot
read the job also cannot adjudicate its meaning or writing quality.

For the first private-job pilot, use customer acknowledgment for payment
and a fixed, bounded refund timeout for unacknowledged jobs. Publish
that policy in the offer before acceptance. This assigns nonpayment risk
to the miner, including when a customer has already consumed a valid
result. Cap pilot job values and record the loss rate. It is an explicit
service tradeoff, not trustless fair exchange. Public protocol evidence
can resolve duplicate or late transactions; it cannot resolve a secret
editorial disagreement.

Benchmark reputation can help customers choose a service without
exposing their own work. It does not prove any particular unknown-input
prediction correct. The customer-directed inspection route in
Section~\ref{customer-directed-evidence-disclosure} can support an agreed
dispute policy. Its effect on payment must be
fixed before job acceptance; a customer-selected opinion alone cannot
redirect escrow. The initial pilot retains acknowledgment and bounded
timeout refunds. Stronger private payment guarantees remain a launch
research question.

\subsection{Revenue and costs}\label{revenue-and-costs}

Customer fees pay for the purchased service. Emissions reward
performance on separate benchmark assignments. Customer volume does not
directly increase emission rewards, because miners can transact with
themselves. Paid successes cannot dilute failures on mandatory protocol
service: the measured availability gate applies in aggregate and
separately to each committed interface and assigned role's mandatory work. Its zero-credit
and recovery-deficit rules are specified in Section~\ref{adversarial-feedback-loop}.
The zero-credit
effect must be validated under the chain's payout lag; it cannot reclaim
already distributed emissions. A fixed quote avoids charging for hidden, unverifiable
internal token or search counts.

Miners pay their own training, inference, private search, and required
candidate Pangram submissions for benchmarks. The benchmark operator or
participating validators fund source-baseline reports from an explicit
round budget; those reports can be reused across competing miners under
the same task/version. Validators also bear A artifact distribution and
benchmark execution, retrieval, hosting, and review costs. Subsidies, if used, must identify their funder and cap
before assignments begin. This design assumes no free API-credit
allocation or automatic on-chain reimbursement scheme.

Assigned B rewrite requests consume reserved provider capacity under
the training-service obligation in
Section~\ref{adversarial-feedback-loop}; the recipient pays no per-call
fee within that budget. Providers bear the compute cost as a condition
of emission eligibility. A miners have requester obligations for those
tickets, with no mandatory A inference endpoint. Authoritative A
benchmark execution uses the validators' separately reserved budget.
B miners may query public A artifacts locally without a protocol quota,
bearing that compute cost themselves. This obligation covers neither
Pangram's external charges nor unlimited hosted search. Bootstrap
reference artifacts and execution need a separately named research
funder. Confidential customer
jobs cannot automatically forward inputs or outputs to Pangram or a
peer miner.
